% Build from this directory: make
% Engine: XeLaTeX. Bibliography: BibTeX + plainnat.
\documentclass[11pt,a4paper,oneside,openany]{book}
\usepackage{fontspec}
\usepackage{xeCJK}
\setmainfont{texgyretermes-regular.otf}[
  BoldFont=texgyretermes-bold.otf,ItalicFont=texgyretermes-italic.otf,
  BoldItalicFont=texgyretermes-bolditalic.otf]
\setsansfont{texgyreheros-regular.otf}[
  BoldFont=texgyreheros-bold.otf,ItalicFont=texgyreheros-italic.otf,
  BoldItalicFont=texgyreheros-bolditalic.otf]
\setmonofont{texgyrecursor-regular.otf}[Scale=0.88,
  BoldFont=texgyrecursor-bold.otf,ItalicFont=texgyrecursor-italic.otf,
  BoldItalicFont=texgyrecursor-bolditalic.otf]
\setCJKmainfont{Noto Serif CJK JP}
\setCJKsansfont{Noto Sans CJK JP}
\setCJKmonofont{Noto Sans Mono CJK JP}
\usepackage[margin=25mm,headheight=15pt]{geometry}
\usepackage{amsmath,amssymb,bm}
\usepackage{booktabs,longtable,array,tabularx}
\usepackage{graphicx,xcolor}
\usepackage{tikz}
\usetikzlibrary{arrows.meta,positioning}
\usepackage[numbers,sort&compress]{natbib}
\usepackage{titlesec}
\usepackage[unicode,colorlinks=true,linkcolor=blue!45!black,
  citecolor=green!35!black,urlcolor=blue!45!black]{hyperref}
\hypersetup{pdftitle={BEACH の物理モデルと数値アルゴリズム},
  pdfsubject={表面帯電、境界要素法、粒子追跡、周期場、準定常シース}}
\titleformat{\chapter}[display]{\sffamily\bfseries\huge}
  {第\thechapter 章}{12pt}{\Huge}
\titleformat{\section}{\sffamily\bfseries\Large}{\thesection}{0.8em}{}
\titleformat{\subsection}{\sffamily\bfseries\large}{\thesubsection}{0.8em}{}
\renewcommand{\contentsname}{目次}
\renewcommand{\listfigurename}{図目次}
\renewcommand{\listtablename}{表目次}
\renewcommand{\figurename}{図}
\renewcommand{\tablename}{表}
\renewcommand{\bibname}{参考文献}
\setcounter{tocdepth}{1}
\setcounter{secnumdepth}{2}
\setlength{\parindent}{1em}
\setlength{\parskip}{0.25em}
\setlength{\emergencystretch}{2em}
\renewcommand{\arraystretch}{1.2}
\newcommand{\vect}[1]{\bm{#1}}
\newcommand{\dd}{\mathop{}\!\mathrm{d}}
\newcommand{\eps}{\epsilon_0}
\newcommand{\kc}{k_{\mathrm c}}
\newcommand{\code}[1]{\texttt{\detokenize{#1}}}
\newcommand{\src}[1]{{\small\nolinkurl{#1}}}
\newcommand{\basisdoc}[1]{\par\smallskip\noindent
  {\small\raggedright\textbf{実装・仕様の根拠：}#1\par}\medskip}
\newcolumntype{Y}{>{\raggedright\arraybackslash}X}

\title{BEACH の物理モデルと数値アルゴリズム\\[0.6em]
  \Large 表面帯電・境界要素法・粒子追跡の技術解説}
\author{技術原稿（著者・所属は確定時に記入）}
\date{2026 年 10 月 6 日\\\small 実装に基づく初稿}

\begin{document}
\pagestyle{plain}
\frontmatter
\maketitle
\chapter*{本書の位置づけ}
\addcontentsline{toc}{chapter}{本書の位置づけ}
本書は BEACH（BEM + Accumulated CHarge）が採用するモデルとアルゴリズムを、
研究背景、支配方程式、離散化、計算手順、適用限界の順にまとめた日本語の技術原稿である。
想定読者は、プラズマと表面の相互作用を研究し、BEACH の計算結果を論文で説明したい研究者である。
設定キーの網羅的な説明は既存の入力リファレンスに委ね、物理量がどの近似を経て計算されるかを重視する。

対象は 2026 年 10 月 5 日に確認したローカル作業ツリーである。
基点の Git commit は \code{57561eb2fe795d84e06e5b6abaf29ff15c1d8070} であり、
確認時点では周期場の長さ尺度と零モードに関する未コミットの変更を含んでいた。
したがって、この原稿を当該 commit のみに対する仕様、あるいは公開リリースの保証と解釈してはならない。
機能の記述には現行 \src{SPEC.md}、対応する説明文書、Fortran 実装を用いた。
インストール済み context plugin の説明よりローカルの新しい実装を優先している。
設定例と粒子源の説明は 2026 年 10 月 6 日に改訂し、境界流入だけの場合は
\code{source_mode} と \code{npcls_per_step} を省略する書き方に揃えた。

本書は新たなシミュレーション結果を報告する研究論文ではない。
性能、実環境との一致、研究上の新規性を実測したとは主張しない。
第\ref{ch:validation}章に、研究ケースごとに必要な検証を示す。
文献は理論の由来と比較対象を示すものであり、BEACH 全体を先行論文と同一のモデルと認定するものではない。

\chapter*{概要}
\addcontentsline{toc}{chapter}{概要}
BEACH は三角形表面要素に蓄積した電荷から静電場を評価し、荷電テスト粒子の前進軌道と
表面への吸収・放出を通じて電荷分布を更新する。各三角形には一定面電荷密度を仮定し、
近傍相互作用を解析的な panel 積分、遠方相互作用を直接和、treecode または Cartesian FMM で扱う。
粒子は予測中点電場を用いた Boris 速度更新と台形位置更新によって進め、軌道線分上の最初の表面衝突または
計算箱境界イベントを処理する。場はバッチ中に固定し、電荷差分をバッチ末尾に確定反映する。
通常の絶縁体電荷蓄積に加え、自由空間での浮遊導体再配分、光電子の照射 ray による放出、
2 軸周期静電場、準定常外部シースとの整合面連成が実装されている。
一方、BEACH 領域内の体積プラズマ空間電荷、誘電分極、表面伝導、外部シースの時間依存輸送は解かない。
本書は、このモデル分担と数値誤差の所在を明確にし、研究利用に必要な収束確認と比較検証を整理する。

\noindent\textbf{キーワード：}表面帯電、境界要素法、テスト粒子、Boris 法、高速多重極法、光電子、準定常シース

\section*{Abstract}
BEACH evaluates electrostatic fields from charge accumulated on triangular surface elements and updates
the charge distribution through forward tracking, absorption, and emission of charged test particles.
Its constant-density triangular sources are evaluated using analytic panel integrals for near interactions
and direct summation, a treecode, or a Cartesian fast multipole method for field evaluation.
Particle positions and velocities are stored at the same time. A predicted midpoint field, Boris velocity
rotation, and trapezoidal position update define each candidate step, followed by the first surface or
box-boundary event. Fields remain fixed during each batch, and charge increments are committed at its end.
Implemented extensions include floating conductors in free space, ray-based photoelectron emission,
two-dimensional periodic electrostatics, and quasistatic coupling to an outer sheath at a matching plane.
The inner domain does not solve volume plasma space charge, dielectric polarization, surface conduction,
or time-dependent outer transport. This technical manuscript documents the model partition, numerical
approximations, and verification procedures needed for research use. It reports no newly generated
simulation results and makes no claim of general experimental validation or measured performance.

\tableofcontents
\mainmatter
\chapter{研究背景と先行研究}
\label{ch:background}

\section{表面電位は粒子の到達を変える}
宇宙空間の表面は、電子・イオンの捕集、光電子放出、二次電子放出などによって帯電する。
表面電位は粒子が到達できる速度領域を変え、その捕集電流が再び帯電状態を変える。
Whipple のレビューは、この電流収支、差動帯電、電位障壁を統一的に整理した基礎文献である\citep{Whipple1981}。
この問題では、総電荷だけでなく、異なる表面部分の電位差が意味を持つ。

単一の浮遊物体を等電位とみなせる場合、電流収支を一つの電位について閉じられる。
絶縁体、多数粒子の集合、粗い地形では、局所的な照射と遮蔽、接近軌道、再吸収先が異なる。
BEACH が三角形要素ごとに電荷を保持するのは、この空間的な差を表現するためである。
ただし絶縁体という指定は、材料内部の分極と伝導まで含むことを意味しない。

\section{月面・無大気天体での動機}
太陽風の電子とイオン、および太陽光で放出される光電子は、無大気天体表面の帯電とシースを作る。
Stubbs らは帯電したダストの静電加速と弾道運動を組み合わせる dynamic fountain model を提案した\citep{Stubbs2006}。
これはダスト移動を考える動機となるが、帯電を求めるだけで浮上量を確定できるわけではない。
付着力、粒子の質量、初期条件を含む別の力学評価が必要になる。

粒子スケールの表面を扱った Zimmerman らの研究は、レゴリス粒子間の電荷差が、
平面シースだけから見積もる電場と異なる局所電場を作り得ることを示した\citep{Zimmerman2016}。
この知見は、幾何と局所電荷を解像する研究の動機となる。
同論文の伝導と絶縁破壊に関する物理まで BEACH に実装されたとは解釈しない。

観測解釈にも複数の機構がある。Horányi らは LADEE の観測した非対称な月周辺ダスト雲を
高速な惑星間ダストの衝突による放出として説明した\citep{Horanyi2015}。
したがって、月面ダストの存在をそのまま静電浮上の証拠と扱うことはできない。
BEACH を使う研究では、帯電機構の検証と、ダスト観測の起源の検証を分ける必要がある。

\section{シースモデルと三次元局所計算}
Zhao らは平面光電子シースの解析解と kinetic PIC 計算を比較し、非単調な Type A と
単調な Type B、Type C の電位構造を扱った\citep{Zhao2020}。
BEACH の Zhao 系 closure は、この平面定常モデルを背景として利用する。
ただし、論文の全領域 PIC を BEACH 内で再現する機能ではない。

Sana と Mishra は、太陽光と太陽風の分布を考慮した単調・非単調シース構造を解析した\citep{SanaMishra2023}。
複数の電位構造が候補となることは、closure の branch を明示して検証する理由になる。
BEACH のエネルギーによる根選択は候補根の比較規則であり、時間依存の安定性解析そのものではない。

三次元の粒子・地形スケールと、より大きい外部シーススケールを両方 kinetic に解くと、
空間解像度、飛行時間、粒子数の要件が同時に増える。
BEACH は局所の三次元表面電荷と軌道に計算を集中し、必要なら外部領域を低次元の準定常応答へ縮約する。
この分担は計算の目的に応じたモデル化であり、PIC と同じ現象を一般に保持する保証ではない。

\section{採用手法の系譜}
三角形上の Green 関数とその勾配の解析積分は、panel 型の場計算の基礎となる。
Graglia は三角形上の線形形状関数を含む積分を整理した\citep{Graglia1993}。
BEACH の通常 source は要素内一定密度の P0 であり、同論文の線形要素をそのまま採用したものではない。

遠方群を階層的にまとめる treecode は Barnes--Hut の考え方に連なる\citep{BarnesHut1986}。
FMM は source の多重極展開と target の局所展開を用いて長距離相互作用を高速化する手法である\citep{GreengardRokhlin1987}。
BEACH は近傍で三角形積分を保ち、遠方に三角形の面積平均 moment を用いる。
このため、三角形を重心の点電荷に置き換えるモデルとは近傍精度も高次 moment も異なる。

2 軸周期系の Ewald 分割については Lindbo と Tornberg の定式化が参考になる\citep{LindboTornberg2012}。
BEACH の通常評価は同論文の FFT ベース solver ではなく、Ewald 参照から生成した遠方 operator と
FMM、さらに独立した物理的零モードを合成する。
粒子更新には Boris 法、三角形交差には Möller--Trumbore 法を利用する\citep{Qin2013,MollerTrumbore1997}。

\basisdoc{\src{SPEC.md}、\src{docs/SurfaceModels.md}、\src{docs/FMMCore.md}、
\src{docs/ZhaoStationaryClosure.md}、\src{docs/MatchingPlaneReference.md}。}

\chapter{モデルの分担と支配方程式}
\label{ch:model}

\section{状態量と計算領域}
表面メッシュを $N$ 個の三角形 $T_i$ で構成し、頂点、面積 $A_i$、重心 $\vect c_i$、
単位法線 $\vect n_i$ を保持する。形状は通常の帯電計算中に固定される。
要素の状態量は総電荷 $q_i$ [C] であり、面電荷密度は
\begin{equation}
  \sigma_i=\frac{q_i}{A_i}\qquad[\mathrm{C\,m^{-2}}]
  \label{eq:p0}
\end{equation}
である。$q_i$ と $\sigma_i$ を混同すると、面積依存の誤りが場と堆積量の両方へ入る。

荷電テスト粒子 $p$ は、位置 $\vect x_p$、速度 $\vect v_p$、実粒子一個の電荷 $q_p$、
質量 $m_p$、macro weight $w_p$ を持つ。軌道の電荷質量比は $q_p/m_p$、
表面へ移す電荷は $q_p w_p$ である。$w_p$ を軌道加速へ重ねて掛けない。
species は電荷・質量だけでなく、速度分布、粒子源、境界作用を指定する。

\section{テスト粒子と表面電荷の自己無撞着性}
粒子の運動方程式は、非相対論的 Lorentz 力
\begin{equation}
 \frac{\dd\vect x_p}{\dd t}=\vect v_p,\qquad
 \frac{\dd\vect v_p}{\dd t}
 =\frac{q_p}{m_p}\left[\vect E(\vect x_p)+\vect v_p\times\vect B_0\right]
 \label{eq:lorentz}
\end{equation}
である。通常経路の $\vect B_0$ は指定した一様外部磁場であり、誘導磁場は求めない。
場は表面電荷と指定した外場から作る。

自由空間の静電場では、表面以外の領域に対して
\begin{equation}
 \nabla^2\phi=0,\qquad \vect E=-\nabla\phi+\vect E_0
\end{equation}
となる。ここで $\phi$ は表面電荷が作る電位であり、一様外場を電位に含める場合は
$-\vect E_0\cdot\vect x$ を一度だけ加える。
追跡中の電子・イオンが空間に作る体積電荷は、この場の右辺に含まれない。

表面電荷の更新と次の粒子軌道は結合するが、体積プラズマを自己無撞着に解く PIC ではない。
外部シース closure を選んでも、BEACH 内部の体積電荷の欠落が自動的に解消されるわけではない。
局所の空間電荷が場を支配するケースでは、その効果を独立した kinetic 計算で調べる。

\section{二つの時間刻み}
粒子の一回の軌道更新幅を $\delta t$、表面電荷更新のバッチ幅を $h$ とする。
$\delta t$ は \code{sim.dt}、$h$ は解決後の \code{sim.batch_duration} に対応する。
\code{batch_duration_step} を使う場合は $h=\delta t\,\code{batch_duration_step}$ である。
バッチ数を $N_b$ とすると、累積帯電時間は受理された幅の和
\begin{equation}
 t_{\mathrm{charge}}=\sum_{n=0}^{N_b-1} h_n
\end{equation}
となる。粒子一個に許す最長追跡時間の目安 $\delta t\,\code{max_step}$ は別の量である。
一つのバッチで複数 species の軌道を順に追っても、その飛行時間を足して帯電時間にはしない。

\section{一バッチの結合手順}
\begin{figure}[htbp]
\centering
\begin{tikzpicture}[node distance=6mm,
 box/.style={draw,rounded corners,align=center,text width=95mm,
   minimum height=9mm,fill=blue!4,font=\small},
 arr/.style={-{Latex[length=2mm]},thick}]
 \node[box] (q) {確定済み表面電荷 $\vect q^{\,n}$};
 \node[box,below=of q] (e) {静電場を構成し、バッチ中は固定};
 \node[box,below=of e] (s) {流入・表面放出の粒子を生成};
 \node[box,below=of s] (p) {Boris 更新と最初の衝突・境界作用};
 \node[box,below=of p] (d) {吸収・放出・必要な closure の電荷差分};
 \node[box,below=of d] (c) {並列集約、候補を検査、受理時に一回だけ反映};
 \node[box,below=of c] (o) {表面電荷 $\vect q^{\,n+1}$ と履歴・電荷台帳を保存};
 \foreach \a/\b in {q/e,e/s,s/p,p/d,d/c,c/o}{\draw[arr] (\a)--(\b);}
\end{tikzpicture}
\caption{基本の表面帯電サイクル。外部シース連成を使う場合は、候補バッチの内側で
同じ初期状態から軌道計算を反復する。}
\label{fig:batch}
\end{figure}
同じバッチ中の粒子は同じ場を参照する。先に吸収された粒子の電荷が、後続粒子へ即時に作用することはない。
この分離によって粒子ごとの並列計算が可能になる一方、場を固定する時間幅の誤差を持つ。
最大 step 数で未解決となった粒子は、解決済みの吸収や escape と区別して集計する。
一般の時間依存粒子 inventory や外部帰還遅延を保持するモデルとして解釈しない。

\basisdoc{\src{src/core/bem_types.f90}、\src{src/particles/bem_particles.f90}、
\src{src/runtime/simulator/bem_simulator_loop.f90}、\src{docs/Algorithms.md}。}

\chapter{三角形境界要素による静電場}
\label{ch:electrostatics}

\section{P0 表面電荷の積分}
自由空間の Coulomb 定数を $\kc=(4\pi\eps)^{-1}$ とすると、
式\eqref{eq:p0}の source が作る電位と電場は
\begin{align}
 \phi(\vect x)&=\kc\sum_{i=1}^{N}\frac{q_i}{A_i}
       \int_{T_i}\frac{1}{\lVert\vect x-\vect y\rVert}\dd S_{\vect y},
       \label{eq:panel_phi}\\
 \vect E_{\mathrm s}(\vect x)&=\kc\sum_{i=1}^{N}\frac{q_i}{A_i}
       \int_{T_i}\frac{\vect x-\vect y}
       {\lVert\vect x-\vect y\rVert^3}\dd S_{\vect y}.
       \label{eq:panel_field}
\end{align}
BEACH の公開 source kernel は \code{triangle_p0} に固定されている。
要素電荷は既知の蓄積量として与えられるため、絶縁体の通常経路では、
境界電位から全 $q_i$ を毎回逆算する連立方程式は解かない。
浮遊導体の等電位化だけが、追加の境界条件を連立方程式へ課す経路である。

P0 は要素ごとの一定密度近似であり、真の表面電荷密度の連続性や細部を完全には表さない。
面の積分を正確に評価しても、三角形形状と密度の離散化誤差は残る。
面上の一定・線形 source に関する解析積分の系譜は Graglia の研究に見られる\citep{Graglia1993}。

\section{辺の対数項と立体角}
解析 kernel の構造を一つの三角形について示す。評価点 $\vect x$ の面からの符号付き高さを
$d$、平面への投影を $\vect x_\parallel$、辺 $e$ の面内外向き法線を $\vect\nu_e$ とする。
辺の長さを $\ell_e$、両端までの距離を $r_{e,1},r_{e,2}$ とおくと、
\begin{equation}
 L_e=\log\frac{r_{e,1}+r_{e,2}+\ell_e}
                   {r_{e,1}+r_{e,2}-\ell_e},\qquad
 d_e=(\vect a_e-\vect x_\parallel)\cdot\vect\nu_e
\end{equation}
を用いて
\begin{align}
 I_\phi&=\sum_e d_eL_e-d\,\Omega,\\
 \vect I_E&=\sum_e\vect\nu_eL_e+\vect n\Omega,\\
 \phi_i&=\kc\frac{q_i}{A_i}I_\phi,\qquad
 \vect E_i=\kc\frac{q_i}{A_i}\vect I_E
\end{align}
を評価する。$\Omega$ は実装の法線規約に従う符号付き立体角であり、
\src{bem_panel_kernel.f90} では三頂点の相対ベクトルの三重積と内積から
\code{atan2} で計算する。
上式は面外での構造を表し、面上では専用の自己項と trace を使う。

重心に点電荷を置く近似と異なり、三角形の広がりを近傍でも保持する。
公開入力に経験的な softening 長を導入して特異性を置き換えるモデルではない。
鋭い辺や頂点の近傍では、電場の特異な振る舞いと数値評価規約を区別する必要がある。

\section{自己電位と片側電場}
要素内の通常の評価点では、面上の電位は連続であり、面積を持つ三角形の自己電位も有限である。
法線電場は表面電荷によって跳ぶ。
\begin{equation}
 \vect E_i^{\pm}
 =\vect E_i^{\mathrm{PV}}\pm\frac{\sigma_i}{2\eps}\vect n_i,
 \qquad
 \vect n_i\cdot(\vect E_i^+-\vect E_i^-)=\frac{\sigma_i}{\eps}.
 \label{eq:jump}
\end{equation}
$\mathrm{PV}$ は両側極限の平均としての principal-value trace である。
どちらを真空・プラズマ側として採用するかは \code{surface_side} と法線の向きで決める。
\code{normal_plus}、\code{normal_minus}、整合した閉曲面の \code{outward_closed} は、
幾何の向きとセットで解釈する。

衝突判定は両側から行われるため、衝突の表裏と電場の片側 trace は同じ概念ではない。
また、面上ちょうどの Direct と FMM には trace 規約の差があるため、
通常の solver 一致試験は面から離した点、または規約を揃えた自己項で行う。

\section{内部の長さ正規化と SI 単位}
代表長 $L$ に対して $\vect x=L\widehat{\vect x}$ と変換した場合、
要素総電荷 $q_i$ を保持したまま、電位は $L^{-1}$、電場は $L^{-2}$ の係数で SI へ戻る。
\begin{equation}
 \phi=\frac{\kc}{L}\widehat\phi,\qquad
 \vect E=\frac{\kc}{L^2}\widehat{\vect E}.
\end{equation}
正規化は数値尺度を整える処理であり、物理的な遮蔽や mesh 細分化を代替しない。
高次 moment は長さの次数が異なるため、周期 operator の生成でも尺度の揃え方が精度に影響する。

\basisdoc{\src{docs/DirectSolver.md}、\src{src/mesh/panel/bem_panel_geometry.f90}、
\src{src/physics/field_solver/panel/bem_panel_kernel.f90}、
\src{src/physics/field_solver/panel/bem_panel_self_terms.f90}。}

\chapter{場の高速化と二軸周期境界}
\label{ch:field}

\section{Direct、treecode、FMM の役割}
$N$ source に対する $M$ 点の Direct 評価は $O(MN)$ の panel 和である。
同じ P0 source のまま遠方近似を導入しないため、小規模の基準計算として利用できる。
全要素重心の電位履歴を Direct で評価すると $O(N^2)$ になり、
粒子追跡とは別に出力が計算時間を支配する可能性がある。

treecode は source の octree を評価点ごとに走査し、十分に遠い node を一つの monopole で置き換える。
FMM は target 側にも局所展開を作り、遠方群の寄与を多くの評価点で共有する。
これらの手法の基礎は Barnes--Hut および Greengard--Rokhlin にある\citep{BarnesHut1986,GreengardRokhlin1987}。
BEACH の \code{auto} は要素数から Direct または FMM を選ぶ。
treecode の通常の対応場は自由空間であり、周期場では Direct/FMM と周期 split の互換性を確認する。

\section{treecode の遠方判定と電荷相殺}
node 内の総電荷、絶対電荷和、電荷中心を
\begin{equation}
 Q_n=\sum_{i\in n}q_i,\quad S_n=\sum_{i\in n}|q_i|,\quad
 \vect c_{Q,n}=\frac{\sum_{i\in n}q_i\vect c_i}{Q_n}
\end{equation}
とする。$Q_n$ が極小のときは幾何中心を用いる。
node 半径 $R_n$ は全三角形頂点を覆い、評価点との距離 $d_n$ に対して
\begin{equation}
 R_n<\theta(d_n-R_n)
\end{equation}
を遠方候補の条件とする。重心だけを用いて半径を過小評価すると、
評価点に近い大きな panel を誤って monopole にまとめる可能性がある。

さらに、正負電荷の相殺による電荷中心の不安定化を避けるため、
\begin{equation}
 |S_n-|Q_n||\le64\epsilon_{\mathrm{mach}}\max(S_n,|Q_n|)
\end{equation}
を満たす node だけを受理する。実質的に異符号電荷が混在する node では子へ降りる。
近傍と leaf は解析 panel 和で評価する。
この規則は相殺時の精度を守る一方、細かく正負電荷が混在する分布では高速化効果を弱める。

\section{Cartesian FMM の展開}
多重指数 $\alpha=(\alpha_x,\alpha_y,\alpha_z)$ を用い、
$|\alpha|=\alpha_x+\alpha_y+\alpha_z$、
$\alpha!=\alpha_x!\alpha_y!\alpha_z!$ とする。
中心 $\vect c$ の panel moment を、階乗を含む規約で
\begin{equation}
 M_\alpha(\vect c)
 =\sum_i\frac{q_i}{A_i\alpha!}
   \int_{T_i}(\vect y-\vect c)^\alpha\dd S_{\vect y}
 \label{eq:panel_moment}
\end{equation}
と定義する。重心点の monomial ではなく、三角形上の厳密な面積平均を用いる。
総電荷 $q_i$ には面積を再び掛けない。

FMM の手順は、leaf の moment を作る P2M、親へ集約する M2M、
遠方 source を target の局所展開へ移す M2L、子へ伝える L2L、
評価点で計算する L2P からなる。近傍は解析 panel の Direct 和を加える。
例えば、中心間変位を $\vect d=\vect c_{\mathrm{child}}-\vect c_{\mathrm{parent}}$ とすると
\begin{equation}
 M_\beta(\vect c_{\mathrm{parent}})
 =\sum_{\mathrm{child}}\sum_{\alpha\le\beta}
 M_\alpha(\vect c_{\mathrm{child}})
 \frac{\vect d^{\beta-\alpha}}{(\beta-\alpha)!}
\end{equation}
で集約する。$G(\vect r)=1/\lVert\vect r\rVert$、
$\vect R=\vect c_t-\vect c_s$ に対して M2L は
\begin{equation}
 L_\alpha(\vect c_t)\mathrel{+}=
 \sum_\beta(-1)^{|\beta|}M_\beta(\vect c_s)
 \partial^{\alpha+\beta}G(\vect R)
\end{equation}
となる。$L_\alpha$ は局所電位の微分係数であり、
\begin{equation}
 \phi_{\mathrm{far}}(\vect x)\simeq
 \kc\sum_{|\alpha|\le p}L_\alpha
 \frac{(\vect x-\vect c_t)^\alpha}{\alpha!},\qquad
 \vect E_{\mathrm{far}}=-\nabla\phi_{\mathrm{far}}
\end{equation}
で評価する。

固定 geometry の木、相互作用 pair、平行移動係数を初期化時に作り、
電荷に依存する moment と local だけをバッチごとに更新する。
現行 simulator の展開次数は $p=4$ 固定である。
低水準 FMM test が次数を変えられることと、\code{beach.toml} に公開次数キーがあることは異なる。
研究ケースでは \code{tree_theta}、leaf size、Direct 基準との差を確認する。
一定次数で良好な木分割を仮定した FMM の計算量の議論を、任意のケースの実測速度として引用しない。

\section{周期場の非零モードと零モード}
\code{periodic2} は $x,y$ 周期、$z$ 非周期の slab を対象とする。
mesh は primitive cell の base triangle を保持し、場と衝突は必要な画像を扱う。
粒子の周期 wrap を指定することと、静電場を無限周期にすることは別である。

cell の横面積を $A_{\mathrm c}=L_xL_y$ とすると、場を
\begin{equation}
 K_{\mathrm s}=
 \underbrace{K_{\mathrm{shell}}+R_{\mathrm{Ewald}}^{\mathrm{full}}
                 -K_0^{\mathrm{sym}}}_{k\ne0\ \text{の backend}}
 +K_0^{\mathrm{physical}}
 \label{eq:periodic_split}
\end{equation}
へ分ける。$K$ は電位と電場の評価作用素を表す。
有限画像だけの設定では、指定した shell の外側を含まない。
\code{field_periodic_far_correction="none"}、および現行の \code{auto} は有限画像モデルである。
無限周期 nonzero mode の production 経路は \code{cached_kneq0} を明示する。
\code{panel_spectral_reference} は小規模な P0 参照経路として用いる。

\section{Ewald 参照から生成する遠方 operator}
Ewald 分割は
\begin{equation}
 \frac1r=\frac{\operatorname{erfc}(a r)}r
            +\frac{\operatorname{erf}(a r)}r
\end{equation}
によって実空間と逆空間の収束を分担する\citep{LindboTornberg2012}。
$a$ は数値計算上の分割 parameter であり、Debye 遮蔽長の逆数ではない。
BEACH の teacher は実・逆空間を有限層で打ち切った Ewald2P 参照である。

初回生成では source の proxy 点と target の check 点を置き、
teacher から有限画像分を引いた残差を再現する operator を正則化 QR で求める。
root multipole を $\vect M_{\mathrm{root}}$ とすると、target anchor $t$ の補正は
\begin{equation}
 \vect L_t^{\mathrm{far}}=\mathsf A_t\vect M_{\mathrm{root}}
\end{equation}
である。電場 fit では決まらない電位の定数係数を、電位残差から別に fit する。
生成した operator を geometry・設定の fingerprint と checksum 付きで保存し、後の run で再利用する。
通常のバッチでは ordinary M2L の後、L2L の前に補正を加えるため、粒子ごとに Ewald 和を再計算しない。

確認した作業ツリーには、高次 moment 行の長さ尺度を揃える生成器修正が含まれる。
それでも、大きな有限画像場と零モードを差し引いて小さな非零場を得る領域では、残差誤差が支配し得る。
cache の再利用と単位尺度の一致だけで、弱い外部場に対する絶対精度を保証できない。

\section{Gauss 則で閉じる物理的零モード}
三角形 $i$ のうち高さ $z$ 以下の面積割合を $F_i(z)$ とおくと、累積電荷は
\begin{equation}
 C(z)=\sum_iq_iF_i(z)
\end{equation}
である。傾いた三角形の $F_i$ は頂点高さの間で区分二次、水平三角形は sheet の跳びになる。
BEACH はこの幾何係数を前計算し、電荷変更時に係数と積分 primitive を更新する。
下側遠方電場 $E_{\mathrm b}$ と gauge $(z_g,\phi_g)$ を指定すると
\begin{align}
 E_{0,z}(z)&=E_{\mathrm b}+\frac{C(z)}{\eps A_{\mathrm c}},
   \label{eq:zero_field}\\
 \phi_0(z)&=\phi_g-E_{\mathrm b}(z-z_g)
   -\frac1{\eps A_{\mathrm c}}\int_{z_g}^{z}C(\zeta)\dd\zeta
\end{align}
となる。全電荷 $Q=\sum_iq_i$ に対する二つの closure は
\begin{equation}
 E_{\mathrm b}=
 \begin{cases}
 -Q/(2\eps A_{\mathrm c}) & \text{対称な上下真空},\\
 0 & \text{下側電束を零に指定}
 \end{cases}
\end{equation}
である。どちらも材料内部の誘電分極や遮蔽を解くモデルではない。
電荷源より上では $C(z)=Q$ を用い、多項式の丸め残差を上部真空へ外挿しない。

非中性 cell の遠方には一定電場と線形電位が残り得る。
また $Q=0$ でも、高さによって正負の電荷が分離すれば $C(z)$ は零とは限らない。
零モードを単に消すと、この高さ分布と Gauss 則を失う。
外部シースを接続する場合も、表面電荷の零モードと外部空間電荷の役割を分け、二重加算を避ける。

\basisdoc{\src{docs/Treecode.md}、\src{docs/FMM.md}、\src{docs/FMMCore.md}、
\src{docs/PeriodicElectrostatics.md}、\src{docs/PeriodicFarCorrection.md}。}

\chapter{粒子積分と幾何イベント}
\label{ch:particles}

\section{同時刻の位置と速度を使う Boris 更新}
入力状態は $(\vect x^n,\vect v^n)$ であり、速度だけ半時刻にずらして保持する状態ではない。
まず予測中点
\begin{equation}
 \vect x_{\mathrm{mid}}=\vect x^n+\frac{\delta t}{2}\vect v^n,
 \qquad \vect E_{\mathrm{mid}}=\vect E(\vect x_{\mathrm{mid}})
\end{equation}
で電場を一回評価する。箱が有効なら、この場の評価点だけを周期軸で wrap、その他の軸で clamp する。
候補軌道の座標は境界の発生順を比較するために物理座標のまま保つ。

速度更新は電場半 step、磁場回転、電場半 step である。
\begin{align}
 \vect v^-&=\vect v^n+\frac qm\vect E_{\mathrm{mid}}\frac{\delta t}{2},\\
 \vect t&=\frac qm\vect B_0\frac{\delta t}{2},\qquad
 \vect s=\frac{2\vect t}{1+\vect t\cdot\vect t},\\
 \vect v'&=\vect v^-+\vect v^-\times\vect t,\\
 \vect v^+&=\vect v^-+\vect v'\times\vect s,\\
 \vect v^{n+1}&=\vect v^++\frac qm\vect E_{\mathrm{mid}}\frac{\delta t}{2},\\
 \vect x^{n+1}&=\vect x^n+\frac{\delta t}{2}
                  (\vect v^n+\vect v^{n+1}).
 \label{eq:trapezoid}
\end{align}
磁場回転は Cayley 型の直交変換となり、電場が零なら速度の大きさを保つ。
ただし数値的な旋回角は $2\arctan(\omega_c\delta t/2)$ であり、
速さを保存しても大きな step の旋回位相まで正確とは限らない。

Boris 法の保存性に関する研究\citep{Qin2013}と、BEACH の場標本化・位置更新は分けて考える。
現行の滑らかな静電場の regression test は位置・速度の二次収束を確認する。
これは衝突、境界作用、バッチ間の場更新を含む全系の二次精度を保証しない。

\section{三角形との最初の交点}
候補軌道の chord を
\begin{equation}
 \vect x(\lambda)=\vect x^n+lambda(\vect x^{n+1}-\vect x^n),
 \qquad 0\le\lambda\le1
\end{equation}
とする。三角形の頂点 $\vect a,\vect b,\vect c$ に対して
\begin{equation}
 \vect a+u(\vect b-\vect a)+v(\vect c-\vect a)=\vect x(\lambda)
\end{equation}
を Möller--Trumbore 法で解き\citep{MollerTrumbore1997}、
$u\ge0$、$v\ge0$、$u+v\le1$、$0\le\lambda\le1$ を満たす交点の最小 $\lambda$ を採用する。
近平行・退化の判定は辺長と線分長に対して尺度を揃える。
法線の符号による back-face 除去は行わず、両側から衝突を検出する。

候補の絞り込みには三角形の AABB を用いる。
小規模 mesh では線形探索、大きい mesh では一様 grid と三次元 DDA によって通過 cell を走査する。
grid の cell に登録された候補三角形の精密交差だけを評価するため、
この grid は静電場の体積 mesh や PIC の電荷割当格子ではない。

\section{計算箱の作用と event の発生順}
mesh hit と box 面通過の fraction を比較し、先に起こるものだけを処理する。
丸め誤差の範囲で同時なら mesh 吸収を優先する。
\begin{center}
\begin{tabularx}{\textwidth}{lY}
\toprule
作用 & 生存粒子への処理 \\
\midrule
\code{open} & escape、または指定した scalar 障壁で return/escape を判定 \\
\code{reflect} & 法線速度を反転し、残り時間を再積分 \\
\code{redistributed_reflect} & 反射に加え、面内の位置を一様に再配置 \\
\code{periodic} & 反対面へ移し、速度を保って残り時間を再積分 \\
\bottomrule
\end{tabularx}
\end{center}
再分布反射は、外部を移動して帰還する粒子の面内位置を縮約する選択肢である。
光電子の上部境界で水平位置を randomize する Zimmerman らの設定に着想を持つ\citep{Zimmerman2016}。
外部軌道の飛行時間や自己無撞着シースを再現する機能とは異なる。

通常 event の速度は chord 方向を使い、その速さを離散仕事
\begin{equation}
 |\vect v_{\mathrm{event}}|^2=|\vect v^n|^2+
 2\frac qm\vect E_{\mathrm{mid}}\cdot
       (\vect x_{\mathrm{event}}-\vect x^n)
\end{equation}
に整合させる。chord の fraction を残り時間の fraction にも使うため、
強い加速・旋回時の真の交差時刻を解いているわけではない。
整合面の $z$ 上端では、coupling 経路が Boris 両端と整合する二次軌道で交差時刻を再評価する。

\section{周期画像と未解決軌道}
周期場の near-image 層数と、衝突で調べる画像範囲は独立に決める。
衝突では物理座標の chord の AABB と canonical mesh の AABB が重なる画像を列挙し、
命中した画像を base 要素へ対応づけて電荷を堆積する。
場評価点を wrap するだけでは、境界の手前で周期画像の表面へ衝突する event を検出できない。

\code{max_step} に達した粒子は未解決として記録する。
不完全な衝突照会を正常な非衝突として扱わない。
一部の周期境界イベントには \code{soft_discard} の選択肢もあるが、失った粒子数と絶対電荷を診断し、
適用範囲を評価する必要がある。未解決量が大きい計算を吸収電流の収束結果とみなさない。

\basisdoc{\src{src/physics/bem_pusher.f90}、\src{src/physics/bem_collision.f90}、
\src{src/runtime/simulator/bem_particle_stepper.f90}、\src{docs/BorisPusher.md}、
\src{docs/ParticleEvents.md}。}

\chapter{粒子流入と光電子放出}
\label{ch:sources}

\section{境界を横切る速度分布}
数密度で規格化した位相空間分布を $f_s(\vect v)$、内向き法線速度を
$v_n=\vect v\cdot\vect n_{\mathrm{in}}$ とすると、境界流入束は
\begin{equation}
 \Gamma_s=\int_{v_n>0}v_n f_s(\vect v)\dd^3v.
 \label{eq:flux}
\end{equation}
境界を横切る粒子は $v_n f_s$ から標本化する。
体積中の Maxwell 分布をそのまま境界速度へ使うと、低速粒子を過剰に注入する。

法線ドリフト $u_n$、一次元熱速度の標準偏差 $s=\sqrt{k_{\mathrm B}T/m}$ を持つ Maxwell 分布では
\begin{equation}
 \Gamma=n\left[
 \frac{s}{\sqrt{2\pi}}\exp\!\left(-\frac{u_n^2}{2s^2}\right)
 +\frac{u_n}{2}\operatorname{erfc}\!\left(-\frac{u_n}{\sqrt2 s}\right)
 \right].
\end{equation}
流入面積 $A_f$ とバッチ幅 $h$ に対し、macro 粒子数の期待値は
\begin{equation}
 \mathbb E[N_{\mathrm{macro}}]=\frac{\Gamma A_f h}{w}.
\end{equation}
整数化の端数を後のバッチへ持ち越すことで、長時間の流入量を表す。

\code{target_macro_particles_per_batch} は標本数から重みを決める指定であり、物理流束の変更ではない。

\section{粒子源と速度 grid}
\code{volume_seed} は指定領域内の標本、\code{plane_source} は内部の軸平行矩形面からの流束を表す。
外部 reservoir は非周期 open 面の \code{boundary_inflow} で与える。
\code{reservoir_face} は deprecated な粒子源であり、新規ケースでは内部面と外部境界を明確に分ける。

境界流入だけの species は、物理パラメータと流入面を指定する次の書き方を基本とする。
\begin{verbatim}
[[particles.species]]
# density, temperature, charge, mass, and macro-particle weight

[particles.species.boundary_inflow]
z_high = "reservoir"
\end{verbatim}
\code{source_mode} と \code{npcls_per_step} は省略する。
省略時の解決値はそれぞれ \code{volume_seed} と \code{0} であり、体積生成は行わない。
体積生成などを使う場合に、その粒子源とパラメータを明示する。

速度 grid には速度三成分と非負の分布値を与える。
\code{phase_space} なら式\eqref{eq:flux}の $v_n$ を掛け、
\code{flux_weighted} なら既に境界流束で重み付けされた分布として扱う。
二重に速度を掛けない。分布形状の規格化と、物理粒子束または電流密度による流入振幅は区別する。

Maxwell 標本化は有限な速度範囲を使い、現行 Gaussian sampling は $6s$ の cutoff を持つ。
通常の熱分布では小さい tail でも、障壁を越える希少な粒子が目的量を支配する場合には確認が必要になる。

\section{無限遠分布の scalar 電位補正}
\code{source_vdf} は境界上で定義した分布を用いる。
\code{infinity_barrier} は上流基準電位 $\phi_\infty$ と境界面平均電位 $\phi_f$ から
\begin{equation}
 \frac12m v_{n,f}^2=\frac12m v_{n,\infty}^2-q(\phi_f-\phi_\infty)
 \label{eq:inflow_energy}
\end{equation}
によって到達条件と法線速度を補正する。$q$ は符号を持つため、電子と正イオンでは応答が逆になる。
必要なエネルギーを持たない上流粒子は到達流束へ含めない。

この補正は面平均による一つの電位障壁の近似である。
箱外の三次元電場、飛行時間、磁化された軌道、途中の非単調電位は自動的に解かれない。
面内電位差が大きいケースでは、面平均流入の誤差を評価する。
外向き open 粒子を単に escape とするか、障壁により反射するかは別の境界作用である。

\section{照射 ray による表面放出}
\code{photo_raycast} は、箱面上の開口から照射 ray を発射し、最初に命中した要素から光電子を放出する。
遮蔽と可視面積が命中率に反映されるため、全表面へ同じ放出量を一律に置くモデルではない。
開口面積 $A_{\mathrm{ap}}$、単位 ray 方向 $\widehat{\vect d}$ に対する投影面積は
\begin{equation}
 A_{\mathrm{proj}}=A_{\mathrm{ap}}
 |\widehat{\vect d}\cdot\vect n_{\mathrm{in}}|.
\end{equation}
全 rank 合計の ray 数 $N_{\mathrm{ray}}$、正の放出電流密度 $J_{\mathrm{emit}}$ に対し、
命中 ray の粒子重みは
\begin{equation}
 w_{\mathrm{hit}}=
 \frac{J_{\mathrm{emit}} A_{\mathrm{proj}}h}{|q|N_{\mathrm{ray}}}.
 \label{eq:ray_weight}
\end{equation}
miss は粒子を生成しない。ray 数を増やすと一粒子の重みが減り、照射積分の統計誤差が小さくなる。
放出時に escape 率を重みに掛けるのではなく、生成後の軌道が再吸収と escape を決める。

法線速度は flux-weighted half-range Maxwell 分布、接線速度は Gaussian 分布から作る。
放出点は入射 ray と反対向きの法線へ微小にずらし、直後の自己衝突を避ける。
確認した実装の変位は $10^{-12}$ m であり、問題の幾何尺度と近接表面間隔に対する相対的大きさも点検する。

放出反作用が有効なら放出元要素 $i$ へ $-q w$、吸収先 $j$ へ $+q w$ を加える。
電子では放出によって正電荷が残り、同じ要素への帰還は相殺、別の要素への帰還は表面間の電荷移送になる。
この移送は粒子の飛行によるもので、材料内の表面伝導ではない。

\basisdoc{\src{docs/ParticleSourcesBoundaries.md}、\src{docs/ReservoirInjection.md}、
\src{docs/PhotoelectronEmission.md}、\src{src/particles/injection/bem_injection_flux.f90}、
\src{src/particles/injection/bem_photoelectron_injection.f90}。}

\chapter{表面電荷更新と時間離散化}
\label{ch:update}

\section{絶縁体への電荷蓄積}
粒子 $p$ が要素 $i$ へ最初に吸収されたとき、その粒子を追跡から除き、
\begin{equation}
 \Delta q_i^{\mathrm{abs}}\mathrel{+}=q_p w_p
\end{equation}
を記録する。放出反作用と選択した closure の補正を含む電荷差分を
thread ごとの配列へ集め、OpenMP の加算と MPI の全体集約の後に一度だけ反映する。
絶縁体の基本更新は
\begin{equation}
 q_i^{n+1}=q_i^n+\Delta q_i^{\mathrm{abs}}
                 +\Delta q_i^{\mathrm{emit}}+\Delta q_i^{\mathrm{closure}}
 \label{eq:charge_update}
\end{equation}
である。一般の表面散乱、二次電子放出、材料内の伝導・拡散はこの吸収モデルに含まれない。

\section{閉じた光電子の統計的帰還 closure}
反射境界によって光電子軌道を閉じても、追跡上限で未帰還が残ることがある。
\code{neutral_return} は、放出粒子の符号付き電荷 $S<0$、解決済み再吸収 $R<0$、
未解決 $U<0$ に対して $S=R+U$ を確認し、帰還 deposit を $S/R$ 倍する。
放出反作用の $-S$ と合わせた表面総電荷増分は零になる。

この closure は未解決粒子の帰還先が、解決済み粒子と同じ分布を持つと仮定する。
長寿命粒子が特定の表面や速度に偏る場合は、この仮定が成立しない可能性がある。
実装は放出があるのに帰還がない状態、実 escape などの不整合を受理せず、
未帰還率 $5\%$ 超も固定の適用範囲外とする。
raw 吸収量と補正後の量を電荷台帳で区別し、倍率と未帰還率を報告する。

\code{fixed_current} は与えた電流 $I$ の目標電荷 $Ih$ に raw 分布を合わせる別の closure である。
光電子では放出と再吸収を別々に規格化する。
大きい二電流の差である小さな正味電流を倍率の分母に用いると、標本誤差を増幅するためである。
この補正は軌道から得た空間分布を使う一方、総電流の振幅は外部モデルから与える。

\section{浮遊導体の等電位化}
自由空間に限り、\code{mesh_id} ごとに浮遊導体の要素電荷を再配分できる。
SI 単位で influence coefficient を
\begin{equation}
 H_{ij}=\frac{\kc}{A_j}\int_{T_j}
       \frac1{\lVert\vect c_i-\vect y\rVert}\dd S_{\vect y}
\end{equation}
と定義する。導体 group $g(i)$ の未知電位を $V_{g(i)}$、
絶縁体電荷と外場の電位を $\phi_i^{\mathrm{fixed}}$ とすると、
\begin{align}
 \sum_{j\in\mathrm{conductor}}H_{ij}q_j-V_{g(i)}
   &=-\phi_i^{\mathrm{fixed}},\\
 \sum_{i\in g}q_i&=Q_g^{\mathrm{before\ relaxation}}
 \label{eq:conductor}
\end{align}
を解く。自己項を含む P0 電位を重心で collocation し、dense system を部分 pivot 付き Gauss 消去で解く。
実装内部では $\kc$ を割った電位規約を用いるが、上式は SI で統一した同じ拘束である。

この再配分は、吸収・放出後の各物体の総電荷を保存する。
物体間に接触電流を流すモデルや、接地電位を固定するモデルではない。
基本の絶縁体計算とは区別して検証し、現行の \code{periodic2} と併用しない。

\section{バッチ幅の安定性}
固定した $\vect q$ に対する粒子応答を十分に平均できると仮定し、
要素ごとの符号付き電流密度を $J_i(\vect q)$ とする。
平均帯電速度 $F_i=A_iJ_i$ に対して、通常の絶縁体更新は概念的に
\begin{equation}
 \vect q^{\,n+1}=\vect q^{\,n}+h\vect F(\vect q^{\,n})+\vect\eta^n
 \label{eq:mean_batch}
\end{equation}
となる。$\vect\eta^n$ は有限標本に由来する電荷誤差である。
滑らかな平均応答の仮定では、場を固定する更新は陽的 Euler 型の帯電離散化と解釈できる。

平均固定点 $\vect q^*$ は $\vect F(\vect q^*)=0$ を満たす。
固定点の式に $h$ が現れなくても、離散更新が同じ点へ安定に収束するとは限らない。
$\mathsf M=\partial\vect F/\partial\vect q|_{\vect q^*}$ に対する線形条件は
\begin{equation}
 \rho(\mathsf I+h\mathsf M)<1
\end{equation}
である。支配固有値が実負で単一の時間尺度 $\tau$ に近似できるときだけ、
$h<2\tau$ が非発散、$h<\tau$ が単調収束の目安となる。
複素固有値、非正規結合、複数固定点にはこの目安をそのまま適用できない。

\section{非零モードの適応幅と監視値}
周期場では候補電荷の変化に対して
\begin{equation}
 \max_i|\Delta\phi_{k\ne0}(\vect c_i;\Delta\vect q)|\le\Phi_{\max}
\end{equation}
を満たすバッチ幅を受理できる。
不合格なら $h,h/2,h/4,\ldots$ を試し、乱数状態と粒子数端数を戻す。
棄却試行の電荷・統計・履歴を確定状態へ混入させない。
この上限は固定した非零場の変化を抑える条件であり、局所打切り誤差推定器や全系の精度保証ではない。
零モード、外部応答、粒子標本の誤差は別に評価する。

電荷変化の監視量は
\begin{equation}
 r_q=\frac{\lVert\vect q^{\,n+1}-\vect q^{\,n}\rVert_2}
 {\max(\lVert\vect q^{\,n+1}\rVert_2,q_{\mathrm{floor}})}
\end{equation}
である。\code{tol_rel} は監視・出力値であり、現行の早期停止条件ではない。
指定した accepted batch 数を処理する。
小さい $h$ や大きな残留電荷によって比が小さくなる場合もあるため、
これだけで正味電流が零、または平衡が達成されたと判断しない。

\basisdoc{\src{docs/SurfaceChargeNumerics.md}、\src{docs/BatchDurationTheory.md}、
\src{src/physics/bem_surface_models_conductor.f90}、
\src{src/runtime/simulator/bem_simulator_loop.f90}。}

\chapter{外部シースの定常・準定常 closure}
\label{ch:sheath}

\section{外部領域を縮約する条件}
BEACH の通常の場は表面電荷に由来し、内側の体積プラズマ空間電荷を含まない。
外部シースが主に平面平均応答を与え、その応答時間が帯電時間より短いと考えられる場合、
外部領域を電位・流束・速度障壁を返す低次元モデルに縮約できる。
これは時間尺度と空間尺度の仮定であり、シースを接続すれば任意の局所 geometry が妥当になるわけではない。

平面・無衝突・非磁化モデルは、電位による速度空間の到達・帰還を使い、
\begin{equation}
 \frac{\dd^2\phi}{\dd z^2}=-\frac{\rho_{\mathrm{plasma}}(\phi)}{\eps}
 \label{eq:outer_poisson}
\end{equation}
を閉じる。単一価 ion、ambient electron、photoelectron の例では
$\rho_{\mathrm{plasma}}=e n_i-e n_e-e n_{\mathrm{pe}}$ である。
分布を到達可能な速度領域で積分して密度を求めるため、単純な Boltzmann 密度だけで全 branch を表せるとは限らない。
Zhao 型モデルの背景は光電子シースの解析・kinetic 比較にある\citep{Zhao2020}。

\section{Type A、B、C と Sagdeev 積分}
上流電位を零とし、表面から外側へ向かう構造を整理すると、
Type A は途中に負の電位極小を持つ非単調な構造、Type B は正の表面から零へ低下する構造、
Type C は負の表面から零へ上昇する構造である。
返ってくる粒子群と通り抜ける粒子群を branch ごとに区別する必要がある。

式\eqref{eq:outer_poisson}に $\dd\phi/\dd z$ を掛け、上流で $E=0$ とすると
\begin{equation}
 \frac{E^2(\phi)}{2}
 =-\frac1\eps\int_0^\phi\rho_{\mathrm{plasma}}(\psi)\dd\psi.
 \label{eq:sagdeev}
\end{equation}
右辺の符号と turning point の条件が、電位 profile の成立性を拘束する。
BEACH の online 経路は密度式、Sagdeev 積分、非線形根、profile の成立性を組み合わせる。
全速度分布の時間発展を解く経路ではない。

現行の Type B ambient electron 密度では、加速後に存在できる粒子の最小速度を反映する。
この下限は 2020 年の Zhao 文献に基づく実装修正として \src{MatchingPlaneReference.md} に説明されている。
未取得の別出版版の式まで同一と確認したとは扱わない。

\section{固定電流の stationary Zhao}
\code{zhao_stationary} は run 開始時に零電流の定常根を一回だけ解く。
ambient electron 吸収 $J_e<0$、ion 吸収 $J_i>0$、光電子放出 $J_{\mathrm{emit}}>0$、
再吸収 $J_{\mathrm{return}}\le0$、escape の表面寄与 $J_{\mathrm{escape}}>0$ とすると
\begin{align}
 J_{\mathrm{return}}&=J_{\mathrm{escape}}-J_{\mathrm{emit}},\\
 J_e+J_i+J_{\mathrm{escape}}&=0,\\
 J_e+J_i+J_{\mathrm{return}}+J_{\mathrm{emit}}&=0.
 \label{eq:stationary_budget}
\end{align}
escape 電流は外へ運ばれる電子電荷に対応する。
これを表面へもう一度 deposit すると、放出反作用と二重計上になる。

求めた branch、基準電位、電位極小、species 別電流 target を run 中固定し、
各バッチの raw hit・return 分布を target に規格化する。
BEACH 内の表面分布と場は更新されるが、零電流根は毎バッチ解き直さない。
光電子なしの stationary 契約では Type C の $J_e+J_i=0$ を解く。

流入は上流の零電位から現在の上端面平均電位へエネルギー保存で写す。
外向きの electron/PE は局所電位と法線運動エネルギーから障壁を越えられるか判定する。
越えられなければ上端で反射し、外部 turning point までの距離と飛行時間は追わない。

\section{整合面での準定常連成}
\code{matching_plane_quasistatic} は箱上端 $z=H$ を外部モデルとの整合面とする。
全表面が $H$ より下にあるとき、内側の平均変位は
\begin{equation}
 D_H=\eps E_{\mathrm b}+\frac{Q}{A_{\mathrm c}}
 \label{eq:matching_d}
\end{equation}
である。外部応答は $D_H$ と粒子が運ぶ moment から、上端電位 $\Phi_H$、
ambient inward flux、各 access potential、光電子 return の障壁を返す。
$(H,\Phi_H)$ を物理的零モードの gauge にすることで、内側の Gauss 則を保ったまま電位を接続する。
外部の field profile を内側に重ねて足す処理ではない。

光電子の外向き粒子束と平均法線エネルギーは、界面通過粒子の重みを使って
\begin{equation}
 \Gamma_{\mathrm{pe}}^+
  =\frac{\sum_{p\in\mathrm{out}}w_p}{A_{\mathrm c}h},\qquad
 \overline{K_n}=\frac{\sum_{p\in\mathrm{out}}w_p m_p v_{n,p}^2/2}
                         {\sum_{p\in\mathrm{out}}w_p}
 \label{eq:matching_moments}
\end{equation}
と表せる。出力・応答では energy を eV に変換する。
online Zhao はこの二 moment を half-Maxwellian に縮約し、
flux-weighted 分布の $\overline{K_n}=k_{\mathrm B}T_{\mathrm{pe}}$ から温度と source 振幅を構成する。
元の高 energy tail や一般の非 Maxwell 分布を保持するわけではない。

\code{table} backend は直積格子の応答表を補間する。
範囲外へ無検証に外挿しない。応答表の grid、上流分布、$H$、単位と符号を研究条件に揃える。
同梱の synthetic CSV は配線確認用であり、物理計算用の応答ではない。
\code{zhao_online} は外部の現在の電束条件から直接根を求めるため、
壁面零電流を run 開始時に一回課す stationary Zhao と拘束が異なる。
特に光電子を省略しても online の branch が必ず Type C になるとは限らない。

\section{固定点反復と陰的零モード}
バッチ開始時の電荷と乱数状態を固定し、外部への feedback $\vect u$ に対して
応答取得と内側の軌道計算を行う写像を $\mathcal T$ と書く。
連成は
\begin{equation}
 \vect u=\mathcal T(\vect u;\vect q^{\,n},h)
\end{equation}
の固定点として求める。例えば緩和係数 $0<\beta\le1$ による反復は
\begin{equation}
 \vect u^{(r+1)}=(1-\beta)\vect u^{(r)}+
                  \beta\mathcal T(\vect u^{(r)};\vect q^{\,n},h)
\end{equation}
である。各 trial は同じ乱数列と粒子数端数から再生される。
独立の乱数で trial ごとに分布を変えると、応答の不一致と標本 noise を分離しにくい。
光電子標本を ambient より先に生成し、ambient 粒子数の変化が光電子の初期標本を変えないようにする。

\code{implicit_zero_mode=true} は、下側電束零の closure で面平均変位だけを
\begin{equation}
 D_H^{n+1}=D_H^n+hJ(D_H^{n+1})
 \label{eq:implicit}
\end{equation}
の後退 Euler 更新にする。table と online はそれぞれの有効範囲で根を bracket して解く。
非零モードの電荷分布、軌道積分、Monte Carlo 誤差まで陰的に解いたことにはならない。

有限のまま固定点反復上限へ達した trial は、残差を記録して警告付きで受理される場合がある。
受理された状態だけが電荷、乱数、history、外部状態を一回更新する。
正常終了と固定点の収束は区別し、\code{matching_plane_history.csv} の残差を確認する。

\section{根の選択と適用限界}
\code{require_unique} は検出した物理根の一意性を要求し、
\code{minimum_energy} は候補のシース電位エネルギーを比較する。
\code{continuation} は明示した Type A と陰的更新のもとで、前の accepted root を seed に使う。
複数構造とエネルギー比較の背景には Sana と Mishra の解析がある\citep{SanaMishra2023}。
有限の multistart は全数学根の isolation ではなく、どの規則も動的安定性や fold 通過を保証しない。

現在の連成は $x,y$ 周期、$z$ open、非磁化、指定外場零などの条件を持つ。
外部の full VDF、粒子 inventory、flight time、遅延 return queue、衝突、過渡シースは解かない。
online Zhao の ambient outward feedback は transparent で、外部での ambient 帰還を解かない。
有限 drift では密度式と流入 VDF の積分が一般に一致しないため、束だけの一致で kinetic 整合性を認定しない。
これらが主要効果となるケースでは、独立の kinetic oracle や PIC を比較対象にする。

\basisdoc{\src{docs/ZhaoStationaryClosure.md}、\src{docs/MatchingPlaneCoupling.md}、
\src{docs/MatchingPlaneReference.md}、\src{docs/OuterKineticOracle.md}、
\src{src/physics/sheath/zhao/bem_matching_plane_zhao_physics.f90}。}

\chapter{検証、収束、再現性}
\label{ch:validation}

\section{実行成功と科学的妥当性}
終了コード零は、設定したアルゴリズムが完了したことを示す。
そのことと、離散化が収束したこと、closure が適用できること、実環境を説明できることは別である。
BEACH の release verification には小規模 fixture の収束記録があるが、
それを研究ケース全体の物理妥当性や性能の証明として引用しない。
本原稿の作成では新しい物理 run や benchmark を実行していない。

\section{誤差を支配する軸から検証する}
\begin{table}[htbp]
\centering\small
\begin{tabularx}{\textwidth}{p{25mm}YY}
\toprule
誤差・近似 & 変える量 & 比較する目的量 \\
\midrule
表面の離散化 & mesh 密度・形状品質 & 表面電位、局所密度、物体別電荷 \\
場の遠方近似 & Direct/FMM、$\theta$、leaf size & 同じ保存電荷に対する $\phi,\vect E$ \\
周期遠方場 & Ewald 層・分割、参照 Fourier & 非零場の絶対誤差、零モードの Gauss 則 \\
軌道と衝突 & $\delta t$、追跡上限 & 命中先、吸収・escape、軌道 energy \\
帯電の更新 & $h$、同じ累積物理時刻 & 電荷履歴、電流、定常分布 \\
有限標本 & 重み・粒子数・ray 数・seed & 平均、分散、低頻度 channel \\
外部 closure & 応答 grid、$H$、明示 branch & $\Phi_H$、流束、return、連成残差 \\
\bottomrule
\end{tabularx}
\caption{研究ケースで切り分ける主な誤差。各軸を同時に変更すると原因を特定しにくい。}
\end{table}

優先する軸は目的量によって異なる。光電子の再吸収先なら ray 標本と衝突幾何、
弱い外部場との比較なら周期 nonzero mode の絶対誤差、
長時間の帯電ならバッチ幅と closure の固定点を先に調べる。
全 parameter に同じ労力を割くより、結論を変え得る近似を重点的に検証する。

\section{独立した oracle を使う基本検証}
無場の直線運動、一様電場の等加速度運動、純磁場の速さ保存と旋回位相を確認する。
滑らかな場では $\delta t,\delta t/2,\delta t/4$ で誤差比を測る。
刻みを小さくしたときは、同じ物理追跡時間を保つよう \code{max_step} も調整する。

場は解析 sheet の jump、有限 panel の独立積分、
同じ P0 mesh の Direct 和、上部真空の Fourier 参照などと比較する。
Direct と FMM の差は FMM 誤差を分離するが、共通の mesh 誤差は検出しない。
零に近い基準場では相対誤差が発散するため、絶対誤差または代表場による規格化も使う。

幾何では一つの三角形への既知の命中、複数要素の最初の交点、
box 面と mesh の発生順、周期画像への衝突を確認する。
境界反射があるケースでは chord 時刻近似も含めて刻み収束を確認する。
密な接触 geometry では表面の重なり、法線、極小間隔の影響を調べる。

\section{電荷台帳と channel 別の収支}
表面の総電荷増分は
\begin{equation}
 Q^{n+1}-Q^n=
  Q_{\mathrm{abs}}+Q_{\mathrm{emit\ reaction}}+Q_{\mathrm{closure}}
 \label{eq:ledger}
\end{equation}
として確認する。導体再配分は各物体内で総電荷を保存するため、右辺へ新たな外部電流を加えない。
開放系では外部から粒子が流入し外へ escape するので、表面だけの $Q$ が保存される必要はない。
閉じた光電子 channel、固定電流 channel、外部境界へ運ばれる電荷を別々に監査する。

raw 吸収、目標電流に対応する電荷、実際に反映した補正電荷を \code{charge_ledger.csv} で区別する。
大きな正負 channel が相殺する場合、正味残差だけでは標本不足を隠す。
倍率、未解決率、各 channel の粒子数と絶対電荷を併せて確認する。

\section{標本誤差と平衡判定}
独立で有限分散の標本に対する Monte Carlo の平均誤差は通常 $N^{-1/2}$ の尺度を持つ。
希少な escape、強い再重み付け、相関した連成反復では単純な標本数だけで有効精度を判断できない。
粒子数や ray 数を増やす比較に加え、複数 seed の反復から主要量の分散を報告する。

バッチ幅を変えると粒子重みや整数化残差も変わる場合がある。
同じ最終バッチ番号より、同じ累積物理時刻で電荷履歴と電流を比べる。
統計的な揺らぎがある平衡では、最後の一点ではなく時間窓の平均と分散を示す。
\code{tol_rel} の小ささだけで平衡を宣言しない。

\section{実装検証と研究成果の再現性}
Fortran の主要 test と Python の出力契約 test は異なる役割を持つ。
前者は場・軌道・収支、後者は設定、結果読込、後処理の契約を確認する。
全体の検証手順は \src{docs/PhysicsReleaseVerification.md} を正本とし、
KUDPC では計算ノード上で実行する。

研究結果には、実際の commit と作業ツリー差分、完全な設定、mesh または OBJ、
応答 CSV、入力速度分布、乱数 seed、MPI/OpenMP 構成、compiler・build 条件を保存する。
周期 cache は cold/warm の区別と generation 条件を残す。
checkpoint 再開では、読込可能だったことと同じ物理条件で継続したことを別に確認する。
並列 reduction と乱数列の違いにより、構成を変えた bitwise 一致は一般に仮定しない。

\basisdoc{\src{docs/ValidationGuide.md}、\src{docs/PhysicsReleaseVerification.md}、
\src{docs/BatchDurationStability.md}、\src{docs/OutputReference.md}、
\src{docs/Execution.md}。}

\chapter{設定例の読み方と研究への接続}
\label{ch:research}

\section{開いている設定例が表す計算}
確認した \src{examples/beach.toml} は、床、小平板、円柱、球を含む複数物体の例である。
床・円柱・球は絶縁体、小平板は浮遊導体を選ぶ。
位置、温度、質量、境界はそのまま特定の月面環境を表す値とは扱わない。

\begin{center}
\small
\begin{tabularx}{\textwidth}{p{57mm}Y}
\toprule
確認した入力 & モデル上の意味 \\
\midrule
\code{dt = 1.4e-8} & 粒子の step 幅 $\delta t=14$ ns \\
\code{batch_duration_step = 200.0} & 帯電更新幅 $h=2.8\,\mu\mathrm{s}$ \\
\code{batch_count = 5} & 固定幅なら累積帯電時間 $14\,\mu\mathrm{s}$ \\
\code{max_step = 500} & 一粒子の最長追跡時間の目安 $7\,\mu\mathrm{s}$ \\
\code{periodic_axes = ["x"]} & 粒子・ray の domain topology の x 周期 \\
\code{field_boundary.mode = "free"} & 静電場は有限 mesh の自由空間場 \\
\code{inflow_model = "infinity_barrier"} & 上流の零電位からの scalar 流入補正 \\
\code{boundary_inflow}\newline\code{z_high = "reservoir"} & 上面全体からの外部 reservoir 流入 \\
\bottomrule
\end{tabularx}
\end{center}
粒子の x 周期と場の自由空間は、独立に選ばれたモデルである。
この例を二軸無限周期のレゴリス層として説明することはできない。
二軸周期場と外部シースを研究する場合は、対応する \code{periodic2} の例から条件を選び直す。

正電荷 species の質量は $9.10938356\times10^{-28}$ kg で、電子質量の $1000$ 倍である。
物理的な陽子をそのまま設定した値ではない。
reduced mass ratio が飛行時間、捕集、シース応答に与える影響を確認する必要がある。
また、例には有効な光電子 species がなく、照射と光電子 feedback を含む月昼側モデルではない。
境界流入だけの species では \code{source_mode} と \code{npcls_per_step} を省略し、
物理パラメータと流入面から意図を読み取れる書き方を採用している。

\section{物体間力と離脱の後処理}
保存済み電荷から対象物体を除く場 $\vect E_{\mathrm{other}}$ を評価すれば、
重心求積による合力と torque は概念的に
\begin{align}
 \vect F&\simeq\sum_{i\in\mathrm{target}}q_i
                 \vect E_{\mathrm{other}}(\vect c_i),\\
 \vect\tau&\simeq\sum_{i\in\mathrm{target}}
       (\vect c_i-\vect x_{\mathrm{ref}})\times
       q_i\vect E_{\mathrm{other}}(\vect c_i)
\end{align}
となる。primary の自己場を除く規約と周期画像の力を保つ規約は、利用する後処理 API で確認する。
BEACH には固定位置の力、移動指標、保存電荷を固定した鉛直離脱経路を調べる後処理がある。

鉛直経路の静電仕事 $W_e(z)=\int_{z_0}^{z}F_z(\zeta)\dd\zeta$ に対して、
重力と付着を含む障壁は
\begin{equation}
 \Delta U(z)=-W_e(z)+mg(z-z_0)+U_{\mathrm{adh}}(z)
\end{equation}
のように評価できる。初期位置で $F_z>mg$ でも、途中の付着・静電障壁を越えられるとは限らない。
保存電荷を動かさずに評価する経路は、移動中の再帯電、周囲粒子の移動、接触の変化を解かない。
その結果は離脱条件の診断であり、ダストの動的放出・飛翔量の予測とは区別する。

\section{論文で明示すべき方法と限界}
BEACH を用いた手法節には、表面 geometry、P0 要素、場境界と solver、
species の分布と重み、$\delta t$ と $h$、追跡上限、表面 model、外部 closure を記載する。
periodic2 なら nonzero mode と物理的 zero mode の closure を、
matching-plane なら応答 backend、branch、$H$、固定点残差を示す。
これらは結果の再現と、モデルの適用性の判断に必要な情報である。

本実装の適用限界には、内側体積空間電荷の省略、材料分極・伝導の省略、
吸収中心の表面相互作用、固定 geometry、バッチ内固定場、外部帰還の即時境界縮約がある。
各研究では、このうち結論に影響する近似を独立の計算や測定と照合する。
例えば内側空間電荷の寄与を PIC と、外部 profile を kinetic oracle と、
粒子放出の初期条件と付着力を実験または別の力学モデルと比較する。

新規性を主張する対象は、採用済みの Boris 法や FMM そのものではなく、
実際に検証した物理的問い、モデル分担、あるいは得られた定量的知見である。
本書の文献調査は選択した基礎・関連研究の確認であり、系統的レビューや網羅的な先行性調査ではない。

\section{データと原稿作成に関する記載}
\textbf{データ利用可能性：}本書の根拠は BEACH のソース、仕様、文書、および参考文献である。
新規計算データは生成していない。実際の研究論文へ展開する際は、設定、mesh、応答表、
収束データの保存先と利用条件を著者が確定する。

\textbf{研究倫理：}本書は計算手法の技術解説であり、人体・動物・個人情報を対象にした研究を含まない。
研究データを追加する際は、その研究に適用される倫理要件を改めて記載する。

\textbf{著者貢献：}著者・所属と CRediT に基づく貢献分担は未確定である。
ソフトウェア開発、検証、原稿の確認と編集の担当を実際の貢献に従って記入する。

\textbf{利益相反・研究費：}情報が提供されていないため未記入である。
「なし」と推定せず、公開・投稿時に著者が確定する。

\textbf{AI 利用：}この初稿は OpenAI Codex による実装読解、文献の書誌確認、文章作成の支援を用いた。
著者による全文の科学的確認や、全引用文献の精読を完了した原稿ではない。
公開・投稿前に、著者が数式、引用箇所、物理的解釈を確認し、対象媒体の方針に沿って開示文を確定する。

\basisdoc{\src{examples/beach.toml}、\src{docs/ObjectForcesDetachment.md}、
\src{docs/PythonPostprocessAPI.md}、\src{SPEC.md}。}

\appendix
\chapter{実装との対応と文献の確認範囲}
\label{ch:sources_appendix}

\section{実装へ戻るための入口}
次のパスは repository root からの相対パスである。
実装の変更時には対応する章の主張と式を更新する。
\begin{longtable}{>{\raggedright\arraybackslash}p{26mm}>{\raggedright\arraybackslash}p{112mm}}
\toprule
内容 & 主な実装・文書 \\
\midrule
\endhead
計算サイクル & \src{src/runtime/simulator/bem_simulator_loop.f90}、\src{docs/Algorithms.md} \\
三角形場 & \src{src/physics/field_solver/panel/bem_panel_kernel.f90}、\src{docs/DirectSolver.md} \\
FMM & \src{src/physics/field_solver/fmm/api/bem_coulomb_fmm_core.f90}、\src{docs/FMMCore.md} \\
周期遠方場 & \src{src/physics/field_solver/fmm/internal/periodic/bem_coulomb_fmm_periodic_root_ops.f90}、\src{docs/PeriodicFarCorrection.md} \\
零モード & \src{src/physics/field_solver/periodic/bem_periodic_zero_mode_plan.f90}、\src{src/physics/field_solver/periodic/bem_periodic_zero_mode_eval.f90} \\
粒子積分 & \src{src/physics/bem_pusher.f90}、\src{src/runtime/simulator/bem_particle_stepper.f90} \\
衝突 & \src{src/physics/bem_collision.f90}、\src{docs/ParticleEvents.md} \\
流入・光電子 & \src{src/particles/injection/bem_injection_flux.f90}、\src{src/particles/injection/bem_photoelectron_injection.f90} \\
導体 & \src{src/physics/bem_surface_models_conductor.f90}、\src{docs/SurfaceChargeNumerics.md} \\
外部シース & \src{src/physics/sheath/zhao/bem_sheath_model_core.f90}、\src{docs/MatchingPlaneReference.md} \\
検証 & \src{docs/ValidationGuide.md}、\src{docs/PhysicsReleaseVerification.md} \\
\bottomrule
\end{longtable}

\section{文献が支える範囲}
\begin{longtable}{>{\raggedright\arraybackslash}p{33mm}>{\raggedright\arraybackslash}p{105mm}}
\toprule
文献 & 本書で用いる範囲 \\
\midrule
\endhead
Whipple (1981) & 表面電位、電流収支、差動帯電の研究背景 \\
Stubbs ら (2006) & 静電的なダスト輸送を考える研究動機。BEACH による飛翔予測の根拠とはしない \\
Horányi ら (2015) & 月ダスト観測に衝突放出の機構があること \\
Zimmerman ら (2016) & 粒子スケール電荷差と境界帰還の背景。材料伝導の実装を意味しない \\
Zhao ら (2020) & 平面光電子シースの Type A/B/C と kinetic 比較の背景 \\
Sana--Mishra (2023) & 単調・非単調構造と候補のエネルギー比較の背景 \\
Graglia (1993) & 三角形上の Green 関数積分の先行研究。公開 source が P1 であるとはしない \\
Barnes--Hut (1986) & source tree による遠方群の近似 \\
Greengard--Rokhlin (1987) & 多重極・局所展開による高速化の理論的系譜 \\
Lindbo--Tornberg (2012) & 2 軸周期 Ewald の分離。BEACH の FFT 実装を意味しない \\
Qin ら (2013) & Boris 法の性質。BEACH 全系の symplectic 性や保存性を主張しない \\
Möller--Trumbore (1997) & ray と三角形の交点を求める方法 \\
\bottomrule
\end{longtable}

DOI、著者、題名、出版年、掲載誌は DOI 登録情報と出版社・研究機関の公開ページで照合した。
内容の確認には公開本文、著者版、abstract を用い、本文未取得の文献について詳細な式の一致を主張しない。
同名の conference 版と journal 版を混同しない。
特に Zhao 文献は DOI \code{10.2514/6.2020-1548} の 2020 年 AIAA 版を引用し、
repository への後年の登録日を出版年として使わない。
\src{references.bib} に DOI と到達先 URL を保存している。

\section{現行仕様と説明文書の優先順位}
モデルの記述を修正するときは、該当する Fortran 実装と現行 \src{SPEC.md} を確認する。
plugin の同梱説明や古い agent guidance には、絶縁体のみを想定した時点の記述が残る場合がある。
実装済み拡張の存在、公開入力から選べる範囲、独立に検証された範囲を別々に記述する。
本書の数式に示した一般的な見方と、公開設定で受理される組合せが異なる場合は、後者を操作の正本とする。

\backmatter
\bibliographystyle{plainnat}
\bibliography{references}
\end{document}
